\section{$\theta_j$ 与旋转频率的关系}

\begin{table}[H]
\centering
\caption{子空间索引与频率特性}
\begin{tabular}{cccc}
\toprule
\textbf{子空间位置} & \textbf{$j$ 值} & \textbf{$\theta_j$} & \textbf{特性} \\
\midrule
左侧（低索引） & 小 & 大（高频） & 快速旋转，编码短距离 \\
右侧（高索引） & 大 & 小（低频） & 缓慢旋转，编码长距离 \\
\bottomrule
\end{tabular}
\end{table}

\begin{knowledgebox}{RoPE 的多尺度编码}
RoPE 不同维度上有不同的旋转频率。将 embedding 向量视为一个行向量时：
\begin{itemize}
  \item 越左侧的 sub-space：频率越高，对相对位置越敏感
  \item 越右侧的 sub-space：频率越低，对相对位置越不敏感
\end{itemize}
这种多尺度特性是 RoPE 成功的关键。
\end{knowledgebox}

